% Résumés : chaque mémoire porte les deux, sa langue d'abord.
\newcommand{\ResumeEN}{%
\chapter*{Abstract}
\addcontentsline{toc}{chapter}{Abstract}
\markboth{Abstract}{Abstract}
\begin{otherlanguage}{english}
The aim of \emph{troncons} is to run, on one ordinary laptop (16\,GB of RAM, an 8\,GB RTX\,4060, a 1\,TB NVMe disk), a language model several times larger than all of its memory, cleanly: weights stored compressed to the extreme on disk, decompressed only for the computation that needs them, at the exact moment it needs them, then released. Read literally, \enquote{without RAM} is impossible; the reading retained here is that \emph{the memory in use no longer depends on the size of the model}, only on a bounded working window.

This first part lays the measured foundations. Every stage between the disk and the compute units was measured on the target machine with tools written for the purpose, each with a control case that was forced to fail: the disk delivers \SI{2.82}{\Go\per\second} (47\,\% of its advertised rate, because the volume is encrypted by BitLocker behind Intel RST VMD), the host-to-device link \SI{12.2}{\Go\per\second}, the video memory \SI{251.3}{\Go\per\second}. From these we prove the roof of the problem, $t \ge \max_i O_i/W_i$, and its corollaries: a crossing factor $r^\star = 4.33$ beyond which compression is wasted if the processor decompresses, and no crossing at all if the graphics card decompresses (architecture~B).

Four levers are then studied and measured. \emph{Extreme quantization}: on the Gaussian source, a shift-register trellis reaches a mean squared error of 0.068 at 2~bits, closing 90\,\% of the gap between the optimal scalar quantizer (0.1175) and the Shannon bound (0.0625); we wrote and measured the first portable trellis decoder (Vulkan, \SI{124}{\Go\per\second} equivalent), and proved a bit-allocation rule by memory stage (with the measured throughputs, a VRAM-resident matrix deserves 2.18 more bits than one read through the link, 3.24 more than one read from disk). \emph{On-demand loading}: what decides everything for a mixture of experts is the cache hit rate; we derive the hit rate implied by any published measurement. \emph{Speculative decoding}: we prove $\mathcal E(\alpha,k)=(1-\alpha^{k+1})/(1-\alpha)$ tokens per read of the weights, and that for fine-grained experts read from disk speculation pays only if $\mathcal E(\alpha,k) > U(k+1)/k_e$, which no independent routing satisfies. \emph{The target models}: 49 open models inventoried from their headers; gpt-oss-120b (63.4\,GB, native MXFP4) is the first \enquote{ultra heavy} target, now downloaded and verified.

\emph{Lossless coding} completes them. The 4-bit codes of the machine's models carry 3.86 bits out of 4, which an exact calculation recovers to the hundredth from any bell-shaped source under a small-group min-max, refuting for these formats the published ``0.6 to 1 bit''. The gain of MXFP4 lies in its scales: with those coded by the median of their row, and decoded exactly, the experts of gpt-oss-120b reach a rate of 1.1026 where zstd takes 1.04. Our interleaved rANS decoder, proved correct, runs at \SI{70}{\Go\per\second} under Vulkan and 82 under CUDA; fused into the matrix--vector product it is 3.7 times slower, so a chunk is decoded once, on arrival in video memory, and gives its whole rate back to the disk.

One equation then contains the four levers: the time of a token of a mixture of experts served by the GPU, stage by stage, under a one-constant locality model. Its proven elasticities say that on this machine the disk holds the maximum (157\,ms per token against 58 for the link and 36 for the GPU, gpt-oss-120b): every bit removed and every byte of throughput gained is worth almost one for one, and faster GPU kernels return under 2\,\%. Three candidate architectures are computed on it. A, exact, serves the 63.4\,GB model with 1.9\,GB of peak RAM. B rests on a new mechanism, the exact refinement of a native 4-bit code into a base and a complement, which costs nothing in entropy (proven) and whose base loses 0.27 bit less than truncation at equal error (measured); counting the complements struck out an earlier $\times 1.42$ gain. The arbiter decided to build A first, in a format that prepares B. First measurements on gpt-oss-20b show a token-to-token expert reuse $\rho = 0.577$, under which speculation would pay from an acceptance rate of 0.42.

The thesis exists in two complete versions, English and French; every figure is regenerated from the committed measurement data.
\end{otherlanguage}
}
\newcommand{\ResumeFR}{%
\chapter*{Résumé}
\addcontentsline{toc}{chapter}{Résumé}
\markboth{Résumé}{Résumé}
\begin{otherlanguage}{french}
Le but de \emph{troncons} est de faire tourner, sur un portable ordinaire (16\,Go de RAM, une RTX\,4060 de 8\,Go, un disque NVMe de 1\,To), un modèle de langue plusieurs fois plus gros que toute sa mémoire, proprement : des poids stockés compressés à l'extrême sur le disque, décompressés seulement pour le calcul qui en a besoin, à l'instant exact où il en a besoin, puis libérés. Au pied de la lettre, « sans RAM » est impossible ; la lecture retenue est que \emph{la mémoire occupée ne dépend plus de la taille du modèle}, seulement d'une fenêtre de travail bornée.

Cette première partie pose les fondations mesurées. Chaque étage entre le disque et les unités de calcul a été mesuré sur la machine cible, avec des outils écrits pour cela, chacun avec un témoin qu'on a forcé à échouer : le disque livre \SI{2.82}{\Go\per\second} (47\,\% de son annonce, parce que le volume est chiffré par BitLocker derrière Intel RST VMD), le lien hôte-carte \SI{12.2}{\Go\per\second}, la mémoire vidéo \SI{251.3}{\Go\per\second}. On en démontre le toit du problème, $t \ge \max_i O_i/W_i$, et ses corollaires : un facteur de croisement $r^\star = 4{,}33$ au-delà duquel la compression est perdue si le processeur décompresse, et aucun croisement si la carte graphique décompresse (architecture~B).

Quatre leviers sont ensuite étudiés et mesurés. \emph{La quantification extrême} : sur la source gaussienne, un treillis à décalage atteint une erreur quadratique de 0,068 à 2~bits et ferme 90\,\% de l'écart entre le scalaire optimal (0,1175) et la borne de Shannon (0,0625) ; nous avons écrit et mesuré le premier décodeur de treillis portable (Vulkan, \SI{124}{\Go\per\second} équivalents), et démontré une règle de partage des bits selon l'étage de mémoire (avec les débits mesurés, une matrice résidente en VRAM mérite 2,18 bits de plus qu'une matrice lue par le lien, 3,24 de plus qu'une matrice lue du disque). \emph{Le chargement à la demande} : pour un mélange d'experts, le taux de succès du cache décide de tout ; on dérive le succès qu'implique toute mesure publiée. \emph{Le décodage spéculatif} : on démontre $\mathcal E(\alpha,k)=(1-\alpha^{k+1})/(1-\alpha)$ jetons par lecture des poids, et que pour des experts fins lus du disque la spéculation ne paie que si $\mathcal E(\alpha,k) > U(k+1)/k_e$, ce qu'aucun routage indépendant ne satisfait. \emph{Les modèles visés} : 49 modèles ouverts inventoriés dans leurs en-têtes ; gpt-oss-120b (63,4\,Go, MXFP4 natif) est la première cible « ultra lourde », téléchargée et vérifiée.

\emph{Le codage sans perte} les complète. Les codes de 4 bits des modèles de la machine portent 3,86 bits sur 4, ce qu'un calcul exact retrouve au centième pour toute source en cloche sous un min-max par petits groupes, ce qui réfute pour ces formats les « 0,6 à 1 bit » publiés. Le gain du MXFP4 est dans ses échelles : avec leurs échelles codées par la médiane de leur ligne, et décodés exactement, les experts de gpt-oss-120b atteignent un taux de 1,1026 là où zstd prend 1,04. Notre décodeur rANS entrelacé, démontré exact, tourne à \SI{70}{\Go\per\second} sous Vulkan et 82 sous CUDA ; fondu au produit matrice-vecteur, il est 3,7 fois plus lent : un tronçon se décode donc une fois, à son arrivée en mémoire vidéo, et rend tout son taux au disque.

Une seule équation contient ensuite les quatre leviers : le temps d'un jeton d'un mélange d'experts servi par la carte, étage par étage, sous un modèle de localité à une constante. Ses élasticités, démontrées, disent que sur cette machine le disque tient le maximum (157\,ms par jeton contre 58 pour le lien et 36 pour la carte, gpt-oss-120b) : chaque bit retiré et chaque octet de débit gagné valent presque un pour un, et des noyaux plus rapides rapportent moins de 2\,\%. Trois architectures candidates y sont calculées. A, exacte, sert le modèle de 63,4\,Go avec 1,9\,Go de RAM de pointe. B repose sur un mécanisme neuf, le raffinement exact d'un code 4 bits natif en une base et un complément, qui ne coûte rien en entropie (démontré) et dont la base perd 0,27 bit de moins que la troncature à erreur égale (mesuré) ; compter les compléments a fait barrer un gain de $\times 1{,}42$ annoncé plus tôt. L'arbitre a décidé de construire A d'abord, dans un format qui prépare B. Les premières mesures sur gpt-oss-20b donnent une réutilisation des experts d'un jeton au suivant $\rho = 0{,}577$, sous laquelle la spéculation paierait dès un taux d'acceptation de 0,42.

Le mémoire existe en deux versions complètes, anglaise et française ; chaque figure se régénère depuis les données de mesure commitées.
\end{otherlanguage}
}
\ifanglais\ResumeEN\ResumeFR\else\ResumeFR\ResumeEN\fi
